\documentclass[10pt,a4paper]{amsart}
\usepackage[margin=19mm]{geometry}
\usepackage{amsmath,amssymb,mathtools}
\usepackage[scr=rsfs,cal=euler]{mathalfa}
\usepackage{xfrac}
\usepackage[unicode,hidelinks,bookmarks=false]{hyperref}
\newcommand{\ie}{i.e.\ }
\newcommand{\cf}{cf.\ }
\newcommand{\resp}{respectively\ }
\newcommand{\Subsection}{\subsection}
\providecommand{\nobreakdash}{\nobreak\hbox{-}}
\setlength{\emergencystretch}{2em}
\multlinegap0pt
\usepackage[shortlabels]{enumitem}
\usepackage{amssymb}
\usepackage{mathtools}
\newtheorem{proposition}{Proposition}
\newtheorem{lemma}{Lemma}
\newcommand{\Bc}{\mathcal{B}}
\title{Log-concavity and unimodality of Hodge numbers of Hilbert schemes of points over a surface}
\author{Anubhab Pahari}
\date{Source: arXiv:2609.04095v1 (CC BY 4.0)}
\begin{document}
\maketitle

\noindent\emph{Source and licence.} Log-concavity and unimodality of Hodge numbers of Hilbert schemes of points over a surface. Version arXiv:2609.04095v1 (CC BY 4.0), distributed by arXiv under the Creative Commons Attribution 4.0 International licence, http://creativecommons.org/licenses/by/4.0/, as recorded in the arXiv licence metadata for this exact version and shown on the abstract page https://arxiv.org/abs/2609.04095v1. Full licence notice: \texttt{licenses/arxiv-2609.04095-CC-BY-4.0.txt}.\par\smallskip

\noindent\textit{Editorial scope.} Counted unit: the article's proposition prop:odd-centres (source0.tex lines 991--994). The article's complete original proof of it is delivered with it (source0.tex lines 996--1191, one \texttt{proof} environment of 196 lines). No mathematics is written, repaired or re-derived: the statement and the proof are the article's own lines, in the article's own notation, and no retained line is substituted or re-flowed.  Retained uncounted as the source-local context the proof needs: lines 352--354; lines 470--472; lines 496--503; lines 505--509; lines 534--549; lines 546--548; lines 584--627; lines 733--737; lines 749--756; lines 770--786; lines 958--967.  Nothing was omitted from any retained span, so the package carries no omission record.  No retained line contains a citation, so the wrapper carries no bibliography and no bibliography entry.  No retained line contains a figure environment, an image-inclusion command or drawing code, so no image is delivered and the package declares no asset dependency.  Editorial formatting only, in addition: an AMS \texttt{amsart} preamble that restates the article's own declarations and macros used by the retained lines, namely the environments \texttt{proposition}, \texttt{lemma} on separate counters, together with the standard packages those lines need.  The retained environments print in this document as Proposition 1, Lemma 1 in order of appearance, whereas the article prints the same items with the numbering of its own full document.  The complete native source of the article as distributed in the arXiv e-print for this exact version is bundled unchanged as \texttt{sources/209343/source0.tex}.
      \begin{equation}\label{eq:universal-tail}
       T_m(W)^2\ge T_{m-1}(W)T_{m+1}(W).
      \end{equation}

\begin{equation}\label{eq:tail-standing-assumptions}
 q\ge2,\qquad g\ge1.
\end{equation}

\begin{equation}\label{eq:w-definition}
 w_u:=
 \begin{cases}
  \displaystyle\binom Nu\alpha_{(m+1-u)/2},
   &u\equiv m+1\pmod2\ \text{and }u\le m+1,\\[5pt]
  0,&\text{otherwise},
 \end{cases}
\end{equation}

\begin{equation}\label{eq:moment-sums}
 Z:=\sum_{u=0}^{N}w_u,\qquad
 V:=\sum_{u=0}^{N}u w_u,\qquad
 \Sigma:=\sum_{u=0}^{N}u(u-1)w_u.
\end{equation}

\begin{proposition}\label{prop:master-identity}
Let \(q,g,W\) be nonnegative integers, and let \(m\ge1\) be an integer.
Then
\begin{equation}\label{eq:master-identity}
\begin{aligned}
 &N^2(N-1)
 \bigl(T_m(W)^2-T_{m-1}(W)T_{m+1}(W)\bigr)\\
 &\hspace{35mm}=(N-1)V^2-NZ\Sigma.
\end{aligned}
\end{equation}
For these fixed \(W\) and \(m\), the inequality
\eqref{eq:universal-tail} is equivalent to
\begin{equation}\label{eq:moment-goal}
 (N-1)V^2\ge NZ\Sigma.
\end{equation}
\end{proposition}

\begin{equation}
 (N-1)V^2\ge NZ\Sigma.
\end{equation}

\begin{lemma}\label{lem:crossing}
Let \(d\ge1\), let \(\omega_0,\ldots,\omega_d>0\), and let
\(\lambda_0,\ldots,\lambda_d\ge0\) be log-concave with interval support
\(\{0,1,\ldots,e\}\), where \(1\le e\le d\).  Put
\[
 \pi_j:=\omega_j\lambda_j,
 \qquad
 M(t):=
 \frac{\sum_{j=0}^d j\omega_jt^j}
      {\sum_{j=0}^d\omega_jt^j}
 \quad(t>0),
\]
\[
 \bar\jmath:=
 \frac{\sum_{j=0}^d j\pi_j}
      {\sum_{j=0}^d\pi_j}.
\]
Then:
\begin{enumerate}[label=\textup{(\alph*)}]
\item \(M\) is strictly increasing on \((0,\infty)\), and there is a
      unique \(t>0\) such that \(M(t)=\bar\jmath\).
\item For this \(t\), choose \(\kappa>0\) such that
      \[
       \widetilde\pi_j:=\kappa\omega_jt^j,
       \qquad
       \sum_j\widetilde\pi_j=\sum_j\pi_j,
      \]
      and put \(\Delta_j:=\pi_j-\widetilde\pi_j\).  Then
      \[
       \sum_j\Delta_j=\sum_jj\Delta_j=0,
      \]
      and \(\{j:\Delta_j>0\}\) is an interval.
\item \(\pi_0\le\widetilde\pi_0\).
\item
      \[
       \sum_jj^2\Delta_j\le0,
       \qquad
       \sum_j\frac{\Delta_j}{j+1}\le0.
      \]
\item If \(\rho>0\) and
      \(\lambda_{j+1}\le\rho\lambda_j\) for \(0\le j<d\), then
      \(t\le\rho\).
\end{enumerate}
\end{lemma}

\begin{equation}\label{eq:parity-polynomials}
 \mathcal E_k(x):=\sum_{j\ge0}\binom{k}{2j}x^{2j},
 \qquad
 \mathcal O_k(x):=\sum_{j\ge0}\binom{k}{2j+1}x^{2j+1}.
\end{equation}

\begin{align}
 \mathcal E_{k-1}(x)^2
 -\mathcal O_{k-2}(x)\mathcal O_k(x)
 &=(1-x^2)^{k-2},\label{eq:parity-first}\\
 \mathcal E_{k-1}(x)\mathcal E_{k+1}(x)
 -\mathcal O_k(x)^2
 &=(1-x^2)^{k-1}.\label{eq:parity-second}
\end{align}

\begin{lemma}\label{lem:reference-sums}
For \(x>0\) and \(L=N-1\),
\begin{align*}
 \sum_{\substack{0\le u\le N\\u\ {\rm odd}}}
 u\binom Nu x^u
 &=Nx\mathcal E_{N-1}(x),\\
 \sum_{\substack{0\le u\le N\\u\ {\rm odd}}}
 u(u-1)\binom Nu x^u
 &=N(N-1)x^2\mathcal O_{N-2}(x),\\
 \sum_{\substack{0\le v\le L\\v\ {\rm odd}}}
 v\binom Lv x^v
 &=Lx\mathcal E_{L-1}(x),\\
 \sum_{\substack{0\le v\le L\\v\ {\rm odd}}}
 \frac1{v+1}\binom Lv x^v
 &=\frac{\mathcal E_{L+1}(x)-1}{(L+1)x}.
\end{align*}
\end{lemma}

\begin{lemma}\label{lem:E-lower-bound}
Let \(L\ge3\), \(0<y<1\), and \(x=\sqrt y\).  Then
\begin{equation}\label{eq:E-lower-bound}
 \mathcal E_{L-1}(x)
 \ge
 (1-y)^{L-3}\bigl(1+(\Bc-2)y\bigr),
 \qquad
 \Bc=\binom L2.
\end{equation}
\end{lemma}

\begin{proposition}\label{prop:odd-centres}
Assume \eqref{eq:tail-standing-assumptions} and \(g\le\Bc\).  Then
\eqref{eq:universal-tail} holds for every odd \(m\ge3\).
\end{proposition}

\begin{proof}
Write \(m=2h-1\) with \(h\ge2\), and fix a nonnegative integer \(W\).
Only even \(u\)
occur in \eqref{eq:w-definition}, and
\[
 w_{2j}=\binom N{2j}\alpha_{h-j}
 \qquad
 \left(0\le j\le\min\left\{h,\left\lfloor\frac N2\right\rfloor\right\}\right).
\]
For every odd \(v=2j+1\) with \(1\le v\le L\), define
\[
 P_v:=(v+1)w_{v+1}
 =
 \begin{cases}
  \displaystyle
  N\binom Lv\alpha_{h-1-j},&0\le j\le h-1,\\[5pt]
  0,&j>h-1.
 \end{cases}
\]
In the remainder of the proof, every sum over \(v\) is over the odd
integers \(1\le v\le L\).  From \eqref{eq:moment-sums},
\begin{equation}\label{eq:odd-dictionary}
 V=\sum_vP_v,\qquad
 \Sigma=\sum_vvP_v,\qquad
 Z=w_0+\sum_v\frac{P_v}{v+1}.
\end{equation}
Since \(P_1=NL\alpha_{h-1}>0\), we have \(\Sigma>0\).

Put
\[
 A:=W+h-1\ge1.
\]
Then
\begin{equation}\label{eq:theta-definition}
 \theta:=\frac{w_0}{P_1}
 =\frac{\alpha_h}{NL\alpha_{h-1}}
 =\frac{A+g}{NL(A+1)},
 \qquad
 w_0=\theta P_1.
\end{equation}
Since \(N-1=L\), condition \eqref{eq:moment-goal} is equivalent to
\begin{equation}\label{eq:odd-goal}
 \theta P_1+\sum_v\frac{P_v}{v+1}
 \le
 \frac{L\left(\sum_vP_v\right)^2}
 {N\sum_vvP_v}.
\end{equation}

Let
\[
 d:=\left\lfloor\frac{L-1}{2}\right\rfloor,
 \qquad
 \omega_j:=N\binom L{2j+1},
\]
\[
 \lambda_j:=
 \begin{cases}
  \alpha_{h-1-j},&0\le j\le h-1,\\
  0,&j>h-1.
 \end{cases}
\]
Then \(P_{2j+1}=\omega_j\lambda_j\).  Put
\[
 e:=\min(h-1,d).
\]
The positive support of \((\lambda_j)_{0\le j\le d}\) is
\(\{0,\ldots,e\}\), and \(e\ge1\).  For \(0\le j<e\),
\[
 \frac{\lambda_{j+1}}{\lambda_j}
 =
 \frac{A-j}{A-j+g-1}.
\]
These quotients are nonincreasing.  Their largest value is
\[
 y_0:=\frac{A}{A+g-1},
 \qquad
 0<y_0\le1.
\]
If \(e<d\), then \(\lambda_{e+1}=0<\lambda_e\), and all subsequent
terms vanish.  Hence \((\lambda_j)\) is log-concave with interval support,
and
\(\lambda_{j+1}\le y_0\lambda_j\) for \(0\le j<d\).

Apply Lemma~\ref{lem:crossing}.  Part~\textup{(e)} gives
\[
 0<t\le y_0.
\]
Put
\[
 y:=t,\qquad x:=\sqrt t,
\]
and put
\[
 \widetilde P_v:=\kappa_{\mathrm o}\binom Lv x^v
 \qquad(v\ {\rm odd}),
 \qquad
 \kappa_{\mathrm o}:=\frac{\kappa N}{x}>0,
\]
so that \(\widetilde P_{2j+1}=\widetilde\pi_j\).
Parts~\textup{(b)},
\textup{(c)}, and \textup{(d)} give
\begin{equation}\label{eq:odd-comparison-data}
\begin{aligned}
 \sum_v\widetilde P_v
 &=\sum_vP_v,\\
 \sum_vv\widetilde P_v
 &=\sum_vvP_v,\\
 P_1&\le\widetilde P_1,\\
 \sum_v\frac{P_v}{v+1}
 &\le
 \sum_v\frac{\widetilde P_v}{v+1}.
\end{aligned}
\end{equation}
Here the last inequality uses \(v+1=2(j+1)\).
Because \(\theta>0\), it is enough to prove
\[
 \theta\widetilde P_1
 +\sum_v\frac{\widetilde P_v}{v+1}
 \le
 \frac{L\left(\sum_v\widetilde P_v\right)^2}
 {N\sum_vv\widetilde P_v}.
\]

By Lemma~\ref{lem:reference-sums} and
\eqref{eq:parity-polynomials},
\[
\begin{aligned}
 \sum_v\widetilde P_v
 &=\kappa_{\mathrm o}\mathcal O_L(x),\\
 \sum_vv\widetilde P_v
 &=\kappa_{\mathrm o}Lx\mathcal E_{L-1}(x),\\
 \widetilde P_1
 &=\kappa_{\mathrm o}Lx,\\
 \sum_v\frac{\widetilde P_v}{v+1}
 &=\kappa_{\mathrm o}
 \frac{\mathcal E_{L+1}(x)-1}{Nx}.
\end{aligned}
\]
After substitution and cancellation, the required inequality becomes
\[
 \theta LNx^2+\mathcal E_{L+1}(x)-1
 \le
 \frac{\mathcal O_L(x)^2}{\mathcal E_{L-1}(x)}.
\]
By \eqref{eq:parity-second},
\[
 \frac{\mathcal O_L(x)^2}{\mathcal E_{L-1}(x)}
 =
 \mathcal E_{L+1}(x)
 -
 \frac{(1-x^2)^{L-1}}{\mathcal E_{L-1}(x)}.
\]
Thus it remains to prove
\begin{equation}\label{eq:odd-star}
 1-\frac{(1-y)^{L-1}}{\mathcal E_{L-1}(x)}
 \ge\theta LNy.
\end{equation}

If \(y=1\), then \(y\le y_0\le1\) gives \(y_0=1\), hence \(g=1\).
Equation~\eqref{eq:theta-definition} gives \(\theta LN=1\), and
\eqref{eq:odd-star} is an equality.

Assume \(0<y<1\).  Lemma~\ref{lem:E-lower-bound} gives
\[
\begin{aligned}
 1-\frac{(1-y)^{L-1}}{\mathcal E_{L-1}(x)}
 &\ge
 1-\frac{(1-y)^2}{1+(\Bc-2)y}\\
 &=\frac{y(\Bc-y)}{1+(\Bc-2)y}.
\end{aligned}
\]
Set
\[
 F(z):=\frac{\Bc-z}{1+(\Bc-2)z}
 \qquad(0\le z\le1).
\]
Since
\[
 F'(z)=-
 \frac{(\Bc-1)^2}
 {\bigl(1+(\Bc-2)z\bigr)^2}<0,
\]
and \(y\le y_0\), we have \(F(y)\ge F(y_0)\).  A direct calculation
using \(y_0=A/(A+g-1)\) gives
\[
 F(y_0)-\frac{A+g}{A+1}
 =
 \frac{(g-1)(\Bc-g)}
 {(A+1)\bigl((\Bc-1)A+g-1\bigr)}
 \ge0.
\]
Since \(\theta LN=(A+g)/(A+1)\), this proves
\eqref{eq:odd-star}.  It follows that \eqref{eq:odd-goal}, and hence
\eqref{eq:moment-goal}, holds.  Proposition~\ref{prop:master-identity}
now proves \eqref{eq:universal-tail} for every odd \(m\ge3\).
\end{proof}

\end{document}
